% !TeX root = ../../Book3.tex
\section{函数芽与解析延拓}
\label{b3:sec:0902}

一个全纯函数的全部导数只描述当前点附近的行为。函数芽把这组局部信息作为一个对象，重叠邻域的恒等关系则提供向外传递的规则。

\subsection{函数芽}
\label{b3:subsec:090201}

\begin{definition}\label{b3:def:09-germ}
在 \(a\) 的邻域全纯的函数按“在更小邻域相同”取等价类，称为在 \(a\) 的\term[hanshu-ya]{函数芽}[germ of a function]，记 \([f]_a\)。
\end{definition}
等价由局部性容易核对。两全纯芽相同，当且仅当在 \(a\) 的 Taylor 系数逐项相同，因为两级数在共同小盘收敛。仅值相同不够：\(1\) 与 \(e^z\) 在0取同值，芽却不同。仅有限多个导数相同也不够：\(0\) 与 \(z^{N+1}\) 的前 \(N\) 阶导数在0均相同。

\ProseParagraphBreak
芽可相加相乘并求导；非零芽有亚纯倒数，只有代表在中心不为零时才有全纯倒数。把函数的某个具体定义域忘掉并不意味着忘掉中心位置，两个不同中心的芽不能直接判作同一个芽。

\ProseParagraphBreak
相同中心处仅有相同值不能确定同一个芽。\(f(z)=z\) 与 \(g(z)=z+z^2\) 在0取相同值、相同一阶导数，但芽不同。所有阶导数相同则由 Taylor 展开保证芽相同。对解析延拓来说，匹配条件是重叠开集上的函数相等，等价于在该重叠连通分支上有一个聚点集的值相等；不能用有限多个插值点取代。

\ProseParagraphBreak
延拓可以超出最初 Taylor 圆盘的收敛半径。\(1/(1+z^2)\) 在0的 Taylor 级数半径为1，但有理表达式在 \(\mathbb C\setminus\{\mathrm{i},-\mathrm{i}\}\) 上全纯，可以沿许多越过单位圆的路径延拓。圆盘收敛半径由离中心最近的奇点限制，不意味着整个圆周都是自然边界。后一节的稀疏幂级数则真有自然边界，两类情况必须用沿路径的局部代表来区分。

\subsection{沿路径延拓}
\label{b3:subsec:090202}

初始芽 \([f]_a\) 沿 \(\gamma\) 的\term[jiexi-yantuo]{解析延拓}[analytic continuation]，是沿路径给一列局部全纯代表，使相邻代表在路径连接处的一个邻域相同，首芽为 \([f]_a\)。紧性允许使用有限子区间和有限小圆盘描述；在每段内，一个代表同时给该段所有位置的芽。

\ProseParagraphBreak
例如初始 \(\log z\) 芽在1取0，可以沿任何去心路径延拓。每个小盘内用 \(1/z\) 的原函数作对数，再以重叠处值匹配常数。绕0一周后，末芽变成初芽加 \(2\pi\mathrm{i}\)。初始 \(1/(1-z)\) 在0的小盘表示式也可沿避开1的路径继续到圆盘外；初始 Taylor 圆的边界不是每一点都真的奇异。

\ProseParagraphBreak
有限链并不要求一个 Taylor 级数走完路径。先在初始盘内取代表，再在相交盘内取与它重叠相等的代表，直到端点。每一次换中心都重算局部收敛盘；连续路径的紧性保证在一项已经存在的延拓上，选出的代表和衔接区可以有限覆盖。这个有限性用于稳定性证明，但不能反过来凭任意盘覆盖宣布代表存在。盘覆盖只给拓扑位置，真正的解析代表仍需函数表达式、方程或其他存在论证。

\ProseParagraphBreak
在有理函数例子中，代表由同一个有理式给出，所以路径只要避开极点就总可延拓。对于根函数，路径必须避开使局部根分支消失的零点，或者允许改用分歧曲面；对于稀疏幂级数，外侧代表可能根本不存在。因此判断“能否沿每条路径延拓”应当先分析阻碍集合，而不是先使用单值性定理来掩盖这项存在条件。

\subsection{唯一性}
\label{b3:subsec:090203}

\begin{lemma}\label{b3:lem:09-continuation-unique}
沿同一条路径的解析延拓若存在，给定初始芽后唯一。
\end{lemma}
\begin{proof}
两延拓在起点附近相同。令 \(E\) 为沿参数已经相同的位置集合；局部代表在邻域相同使 \(E\) 开。若相同位置趋于某个参数，取该参数附近两套代表的小共同圆盘，相同位置的芽使两代表在其中有开集相同，恒等定理使中心芽也相同，所以 \(E\) 闭。区间连通且 \(E\) 非空，故是全部区间。
\end{proof}
\begin{lemma}[有限链稳定性]\label{b3:lem:09-continuation-stability}
若初始芽可沿 \(\gamma\) 延拓，则充分接近 \(\gamma\) 的路径也可使用同一组有限局部代表延拓；若两路径起终点相同，末芽相同。
\end{lemma}
\begin{proof}
把路径分成有限小段，使每段的像紧含于一个代表圆盘，且相邻接点紧含于两个代表相同的共同小邻域。有限个正的距离下界给一个共同 \(\varepsilon>0\)。与原路径逐点距离小于 \(\varepsilon\) 的新路径仍落在相同段盘中，接点也在相同共同邻域，所以可逐段用同一代表。终点固定时最终使用同一末代表，故芽相同。
\end{proof}
稳定性只对已有延拓的固定紧路径给一个邻域。它不说明越过某个奇点的路径也可继续，或任意长路径都有统一延拓半径。

\subsection{最大解析延拓}
\label{b3:subsec:090204}

从一个初始芽出发，把沿允许路径能到达的全部芽收集起来。以局部代表 \(f\) 的图式 \(z\mapsto[f]_z\) 作为坐标片，得到一个投影到基区域的局部双全纯空间；函数值映射在这些坐标片上就是 \(f(z)\)。这称为该芽的最大解析延拓空间，最大性指任何其他沿同一区域的延拓都通过芽唯一地映入它。

\ProseParagraphBreak
不同基点可用不交基圆盘分开；同一基点的不同芽也可分开，因为若两个小盘代表在任一共同开子集相同，恒等定理迫使中心芽相同。因此芽空间是 Hausdorff 的。通过有理中心和有理半径的有限延拓链可选出可数坐标基，所以它是一张通常意义的 Riemann 曲面。其投影局部单射；要把代数函数的分歧点补入，则需另外完成曲面，投影将不再是局部单射。

\ProseParagraphBreak
对 \(\sqrt z\)，在去心平面上芽空间是二层覆盖，可用 \(w\in\mathbb C^*\)、\(z=w^2\) 表示。补入 \(w=0\) 后得到整张平面，但 \(w\mapsto w^2\) 在0分歧。补点后的对象有用，却不应与“基变量 \(z\) 的全纯芽延拓”混为一谈。

\ProseParagraphBreak
最大延拓依赖所允许的对象。若只允许全纯芽，\(1/z\) 的芽曲面不能在0上方加一个全纯点；若允许球面值的亚纯芽，就可以在0补一个取无穷远值的点。类似地，\(\sqrt z\) 在 \(\mathbb C^*\) 上的全纯芽空间投影无分歧，而用 \(z=t^2\)、\(F=t\) 补上 \(t=0\) 后，得到延伸到平面底层的分歧曲面。底层区域扩大、允许极点或允许分歧，是三个不同操作，不能都含糊称为“补一个点”。

\ProseParagraphBreak
芽空间的局部图由一个代表 \(f\) 在小盘中的全部芽 \(\{f_z\}\) 构成，投影在这一图上是一一对应的小盘坐标。若两个不同芽在同一中心，选共同小盘上的代表；恒等定理使它们要么处处同芽，要么在该小盘每个中心都不同。因此两张小盘图可分开，给 Hausdorff 性。可达性还使整个空间连通：沿一条延拓路径把有限链逐段拼成芽空间中的路径即可。局部双全纯投影只有小图信息，并不自动满足覆盖映射要求的完整均匀层分解；后面将再次区分。
